\subsection{Core-side primitive annuli with passive geometric transport}
\label{marked-cut:core-annuli}

A bundle whose whole base boundary is attached need not supply the dual disks
used in Section~\ref{marked-cut:primitive-annuli}. In five of the remaining
source examples, the required disks instead exist in the simplified core.
The new certificate authenticates the annular midpoint curves and carries
their unsigned intersection counts through the actual core-cancellation
schedule. A surviving cocore that meets one pending annulus once and every
other pending annulus zero times is a \emph{private cocore} for that annulus.
A complete set of these cocores certifies a simultaneous primitive family
on the core side.

\paragraph{Source curves precede transport.}
Every full attaching circle on a nondisk base is traced using its signed
port word. A zero-dimensional port stays in its source zero-handle. A
one-dimensional port must be a face-sector side of an edge-cap polygon;
the ordered tetrahedron/face flags authenticate its direction through the
source one-handle. Consecutive endpoint owners must form a closed core path.
The curve is the midpoint section of the interval fibres. Its embeddedness
and disjointness from the real two-handle attaching curves follow from the
source product-intersection theorem, including the actual disk bundles
already absorbed as two-handles.

The extracted chart routines are shared with disk-base absorption, but the
producer and independent checker retain separate implementations. The
checker reconstructs endpoint flags and paths from authenticated source
geometry. An arbitrary group word or a claimed incidence row cannot supply
this prerequisite.

\begin{proposition}[Passive geometric pivot transport]
\label{marked-cut:passive-pivot}
Let the authenticated core have a successful geometric cancellation schedule,
and let $\beta_1,\ldots,\beta_q$ be pairwise disjoint annular midpoint curves
outside the real two-handle attaching regions. If a real pivot eliminates
row $r$ and cocore column $g$, then the passive counts satisfy
\[
 m'_{ih}=m_{ih}+m_{ig}n_{rh}\qquad(h\ne g),
\]
where $n$ is the current real attaching-curve incidence matrix. This formula
holds simultaneously for all passive curves. Saturation at two preserves
exactly which final entries are zero or one.
\end{proposition}
\begin{proof}
Apply the disk slides in Lemma~\ref{marked-cut:pivot}. Every sliding band
follows the real attaching curve $\alpha_r$ and has interior disjoint from
all other real and passive curves. Each slide of $D_h$ over $D_g$ adds
exactly $m_{ig}$ intersections with $\beta_i$. There are $n_{rh}$ such
slides. After the unique-intersection cancellation, $\beta_i$ remains a
boundary curve and has the displayed intersection counts with the surviving
complete disk system. The same bands work for every passive curve because
the curves are pairwise disjoint from the band interiors.

The initial counts are actual unsigned source crossings; opposite traversal
signs do not cancel. Every subsequent expression uses only addition and
multiplication of nonnegative integers. The map $x\mapsto\min(2,x)$ preserves
both operations in the saturated semiring, so it preserves the final zero
and singleton predicates even if the unsaturated counts become very large.
\end{proof}

\begin{proposition}[Private cocores certify a simultaneous family]
\label{marked-cut:core-private-family}
Suppose that, after the real core cancellations, every pending annulus $i$
has a surviving column $h_i$ with $m_{i h_i}=1$ and $m_{j h_i}=0$ for all
$j\ne i$. Then the annuli form a simultaneous primitive family in their
respective core handlebodies. The columns $h_i$ are necessarily distinct.
\end{proposition}
\begin{proof}
The surviving cocores belong to one embedded complete disk system, so they
are pairwise disjoint. The stated counts say that each designated disk meets
its midpoint curve once and the other midpoint curves zero times. Replace
each attaching annulus by a sufficiently thin regular neighborhood of its
midpoint curve, using a boundary isotopy. Its intersection with its designated
disk is then one spanning arc, and its intersections with the other designated
disks are empty. These isotopies are simultaneous because the annuli are
disjoint. Finite proper disk-attachment regions, also disjoint from the annuli,
can be avoided by a small simultaneous boundary isotopy of the cocore disks.
Thus the primitive family is relative to those remaining disk interfaces.
A column cannot serve two curves, since it would have to contain both a
singleton and a zero at one of their rows.
\end{proof}

\paragraph{Whole-cut gluing and ownership.}
Use the annulus-gluing one/two-handle description in
Proposition~\ref{marked-cut:primitive-census}, with the private handle now on
the core side. Each new two-handle cancels its private core one-handle, while
the connecting handle remains. The disks are simultaneous, so one gluing
does not invalidate another. Core components inherit orientability from the
source cut. For each connected attachment graph component,
\[
 g=\sum_{v\in V}(g_v-a_v)+|E|-|V|+1,
\]
where $a_v$ now counts the private annuli owned by core vertex $v$, and is
zero at bundle vertices. Here $E$ includes both proper disk interfaces and
all remaining nondisk annuli. Each selected cocore must belong to the same
verified core component as its source curve; the whole-cut graph checker
also checks that ownership. The original source Euler identity is unchanged.

The current producer first tries the complete bundle-side rule. If that fails,
it defers every nondisk circle, absorbs only genuine disk-base components,
and requires a successful core cancellation. It then attempts the private
column test on \emph{all} pending nondisk circles simultaneously. It does not
mix a partial bundle-side family with a partial core-side family. This avoids
assuming compatibility between two independently chosen disk systems. A
failed core schedule or a passive row without a private singleton remains
inconclusive.

\paragraph{Independent replay.}
The new whole-cut schema is \path{normal-cut-handlebodies-v4}. The annular
record contains each source attachment and base owner, its directed source
crossings, its final saturated row on the surviving columns, and its private
column and core component. The enclosing verifier authenticates the source
parallelity and boundary proofs, checks actual disk absorption, and verifies
the augmented core cancellation before interpreting this record.

The producer transports sparse rows. The independent module
\path{core_primitive_annuli_verify} reconstructs both real and passive rows,
then performs a dense replay of every supplied real pivot. It recomputes the
canonical least private column and compares every recorded field. It imports
no producer. Only after all these checks does the final graph computation
deduct the corresponding core genera. Schemas v1--v3 remain replayable;
a downgrade to an earlier schema cannot silently omit the passive proof.

This is an authenticated transport of the \emph{midpoint intersection
predicates} needed for a terminal genus census. It does not construct the
possibly long disk slides or encode every later boundary-pattern strand.
A complete hierarchy still needs more geometric information.

\paragraph{Cost and incomplete search.}
There are $O(T)$ real handles and $q\le P=O(T)$ passive curves. At most
$O(T)$ real pivots update an $(O(T)+q)$-by-$O(T)$ matrix. Dense replay and
the private-column check take $O(T^3)$ operations on saturated entries.
The added final rows use $O(PT)$ entries with $O(\log(T+P))$-bit labels;
the directed source paths together use $O(P)$ crossings. Normal multiplicity
is never expanded. These bounds add polynomial work and certificate size to
the existing source queries.

The success criterion depends on the supplied core schedule and disk system.
The current greedy core producer is incomplete and label-dependent. In a
fixed five-relabel test, three sources certify the same census and two stop
at the real core-cancellation stage before passive transport. A different
cocore system could also expose primitivity when the current private-column
test fails. Neither a failed schedule nor absence of a private column is a
negative topological certificate. Global useful-surface coverage and the
quasi-polynomial hierarchy bound remain unresolved.

\paragraph{Source audit and preserved certificates.}
The refreshed audit certifies 83 of the 106 supplied cuts. All 83 source
certificates replay independently. Regina constructs the actual cut and
recognizes all components in 99 cases, agreeing exactly with each of the
76 native successes within the 4,096-normal-disk expansion cap. Seven large
successes use source replay without literal expansion, and no oracle times out.
Data and runtime pins are retained in \path{data/core-annuli-audit.json}.

All 78 previously accepted inputs remain accepted, and their canonical
certificate hashes are unchanged. The five new v4 certificates are exactly
the private-column candidates identified by the earlier passive probe.
Each of those source cuts consists of one solid torus and one genus-two
handlebody, in agreement with the explicit oracle.

The 23 remaining refusals comprise two annular cases without a complete
private-column family, eleven failed core cancellations (including three
square raw obstructions), and ten cores with three-handles. Four cases that
were previously classified first by their unsupported annuli now reach and
fail the real core-cancellation stage. No earlier success is lost.

Thirty focused tests pass. They cover the five source families, hostile
source crossings and final rows, incorrect cocore owners, omitted pivots,
schema downgrades, disabled producers, hexadecimal transport, budget caps,
large binary multiplicity and earlier certificate versions. Synthetic
conditional-chart tests compare the passive update $(2,1,1)\mapsto(0,3,1)$
with its saturation $(0,2,1)$, including opposite signs on the repeated
crossing. They reject two annuli competing for the same alleged private
column. Those algebraic test charts are not asserted to be source cuts.
The five-relabel source test explicitly records three accepted censuses and
two conservative core-stage refusals.

\paragraph{Full regression.}
The complete Regina-enabled suite passes all 1,671 tests with no skips in
523.669 seconds. Its 13,846 expanded-sheet topology comparisons and 196,712
independent marked-cover comparisons also pass. The log is retained in
\path{data/core-annuli-tests.txt}; the runtime source pins match the 106-case
audit. These checks validate the supplied-cut implementation and its refusal
contracts, not completeness of useful-surface discovery.

